% !TeX root = ./main.tex

% --------------------------------------------------
% 資訊設定（Information Configs）
% --------------------------------------------------

\ntusetup{
  university*   = {National Taiwan University},
  university    = {國立臺灣大學},
  college       = {工學院},
  college*      = {College of Engineering},
  institute     = {工程科學與海洋工程研究所},
  institute*    = {Institute of Engineering Science and Ocean Engineering},
  title         = {應用知識圖譜注意力網路於食譜推薦之效能評估與可解釋性研究},
  title*        = {Performance Evaluation and Explainability Study of Knowledge Graph Attention Network for Recipe Recommendation},
  author        = {廖健合},
  author*       = {Jian-He Liao},
  ID            = {R12525066},
  advisor       = {張瑞益},
  advisor*      = {RaY-I Chang},
  % date          = {2026-04-30},         % 若註解掉，則預設為當天
  % oral-date     = {2026-07-30},         % 若註解掉，則預設為當天
  % DOI           = {10.5566/NTU2026XXXXX},
  keywords      = {知識圖譜, 注意力機制, 食譜推薦, 深度學習, 協同過濾, 可解釋性},
  keywords*     = {Knowledge Graph, Attention Mechanism, Recipe Recommendation, Deep Learning, Collaborative Filtering, Explainability},
}

% --------------------------------------------------
% 加載套件（Include Packages）
% --------------------------------------------------

\usepackage[numbers,sort&compress]{natbib} % 參考文獻（數字格式，與 abbrv 相容）
\usepackage{amsmath, amsthm, amssymb}   % 數學環境
\usepackage{ulem, CJKulem}              % 下劃線、雙下劃線與波浪紋效果
\usepackage{booktabs}                   % 改善表格設置
\usepackage{multirow}                   % 合併儲存格
\usepackage{diagbox}                    % 插入表格反斜線
\usepackage{array}                      % 調整表格高度
\usepackage{longtable}                  % 支援跨頁長表格
\usepackage{paralist}                   % 列表環境
\usepackage{subcaption}                 % 子圖
\usepackage{enumitem}                   % 進階列表
\usepackage{bm}                         % 粗體數學符號
\usepackage{tabularx}                   % 自適應寬度表格
\usepackage{float}                      % 強制定位表格 [H]
\usepackage{xcolor}                     % 顏色支援
\usepackage{pifont}                     % 特殊符號 (ding)
\usepackage{standalone}                 % 獨立輸出
\usepackage{tikz}                       % 繪圖
\usetikzlibrary{positioning, arrows.meta, shapes.geometric, fit, backgrounds, shadows}
% --------------------------------------------------
% 套件設定（Packages Settings）
% --------------------------------------------------